
%======================================================================
\section{Leverage and Influence}\label{sec:leverage}
%======================================================================

\subsection{Leverage as geometric distance}

Each fitted value $\hat{y}_i = \sum_j h_{ij}y_j$ is a weighted
combination of all responses. The diagonal element $h_{ii}$
measures how much observation $i$ influences \emph{its own} prediction.

\begin{definition}
The \textbf{leverage} of observation $i$ is $h_{ii}$, the $i$-th diagonal
element of $\bH$. By Proposition~\ref{prop:hat}(iv),
$0 \le h_{ii} \le 1$ and $\sum_i h_{ii} = p$.
The conventional threshold for ``high leverage'' is $h_{ii} > 2p/n$.
\end{definition}

\begin{remark}[The bound $h_{ii}\ge 1/n$ when an intercept is present]\label{rem:hii-intercept}
When the model includes an intercept (i.e., a column of ones $\bm{1}\in\Col(\bX)$),
the leverage satisfies the stronger bound $h_{ii}\ge 1/n$.
To see why: the hat matrix projects onto $\Col(\bX)$, which contains $\bm{1}$.
In particular, $\bH\bm{1} = \bm{1}$, so the $i$-th row of $\bH$ sums to 1:
$\sum_{j=1}^n h_{ij} = 1$.
This can be proved using the Cauchy--Schwarz inequality applied to the
rows of~$\bH$; the proof is in Appendix~\ref{app:proofs}.

Without an intercept, it is possible to have $h_{ii} = 0$ (for an observation
$\bm{x}_i$ orthogonal to every column of $\bX$), so the general bound is
$0 \le h_{ii} \le 1$.
\end{remark}

Geometrically, $h_{ii}$ measures how far the $i$-th predictor vector
$\bm{x}_i^\top$ lies from the centroid of the predictor space.
Observations in unusual predictor combinations have long ``arms''
in the tug-of-war that determines the projection.

\subsection{Cook's distance: leverage $\times$ surprise}

High leverage means an observation \emph{can} distort the projection.
Whether it \emph{does} depends on whether it also has a large residual.
Cook's distance combines both pieces of information into a single measure of \emph{influence}.

\medskip
\noindent\textbf{The thought experiment behind Cook's distance.}
Imagine deleting observation $i$ from the dataset and re-running the entire regression.
How much do the fitted values change?
If they barely move, observation $i$ is not influential.
If they change dramatically, observation $i$ is pulling the regression toward itself.

Cook's distance measures exactly this: it is proportional to the squared distance between the fitted values from the full regression and the fitted values with observation $i$ removed:
\[
  D_i \;\propto\; \norm{\yhat - \yhat_{(i)}}^2,
\]
where $\yhat_{(i)}$ denotes the fitted values from the regression that excludes observation $i$.
The remarkable fact is that you do not actually need to refit the regression $n$ times.
Using the hat matrix, Cook's distance can be computed in closed form.

\begin{definition}
\textbf{Cook's distance} for observation $i$ is
\[
  D_i = \frac{e_i^2}{p\cdot\text{MSE}}
        \cdot \frac{h_{ii}}{(1 - h_{ii})^2},
\]
where $\text{MSE} = \norm{\be}^2/(n-p)$ is the mean squared error.
\end{definition}

\noindent\textbf{Reading the formula.}
Cook's distance is a product of two factors:
\begin{enumerate}[nosep]
  \item $\frac{e_i^2}{p\cdot\text{MSE}}$ --- this is the \textbf{surprise factor}.
    It measures how large the residual of observation $i$ is, relative to the overall noise level.
    A large residual means the observation does not fit the current model well; it ``wants'' the regression plane to tilt toward it.
  \item $\frac{h_{ii}}{(1 - h_{ii})^2}$ --- this is the \textbf{leverage factor}.
    It measures how much ``mechanical advantage'' the observation has.
    An observation at an extreme predictor value ($h_{ii}$ close to 1) can exert enormous force on the regression; an observation near the center of the predictor cloud ($h_{ii}$ close to $1/n$) has little leverage, no matter how unusual its $y$-value.
\end{enumerate}

\medskip
\noindent\textbf{An analogy.}
Think of the regression line as a seesaw (lever) balanced on the centroid of the data.
\emph{Leverage} ($h_{ii}$) is how far from the fulcrum observation $i$ sits---its position along the seesaw.
The \emph{residual} ($e_i$) is the force it pushes with---how far above or below the seesaw it is.
Cook's distance is the resulting \emph{torque}: force times distance from the fulcrum.
A person sitting at the end of the seesaw (high leverage) and pushing hard (large residual) tilts the seesaw a lot. A person sitting near the fulcrum (low leverage), no matter how hard they push, barely tilts it.

\medskip
\noindent\textbf{When to worry.}
A common rule of thumb is to investigate observations with $D_i > 4/n$ or $D_i > 1$.
But the real message is qualitative: look at observations with large Cook's distance
and ask whether they represent data errors, unusual subpopulations, or genuine outliers
that should influence the model.

\begin{keyinsight}
\textbf{Influence = leverage $\times$ surprise.}
An observation needs \emph{both} a long arm ($h_{ii}$ large) \emph{and}
a strong pull ($|e_i|$ large) to actually move the regression.
A high-leverage observation that fits the trend is harmless
(large $h_{ii}$, small $e_i$, so $D_i$ is small).
A central observation with a large residual has limited impact
(small $h_{ii}$, large $e_i$, so $D_i$ is moderate).
The dangerous case is an observation that is \emph{both} extreme in its
predictor values \emph{and} does not follow the pattern of the other data.
\end{keyinsight}

% --- Figure 4: Multicollinearity (placed here as visual bridge to Section 7) ---

\begin{figure}[ht]
\centering
\begin{tikzpicture}[>=Stealth, scale=1.0]

  % ============== LEFT: Well-conditioned ==============
  \begin{scope}[shift={(-4.8,0)}]
    \node[secondary, font=\small\bfseries] at (0,3.6) {Well-conditioned};

    % Column vectors at ~60 degrees
    \draw[->, colspace, very thick] (0,0) -- (2.6,0.5)
      node[right, font=\footnotesize\bfseries] {$\bm{x}_1$};
    \draw[->, colspace!70!black, very thick] (0,0) -- (0.8,2.5)
      node[above, font=\footnotesize\bfseries] {$\bm{x}_2$};

    % Projection yhat
    \coordinate (yhat) at (1.8,1.6);
    \draw[->, projection, thick] (0,0) -- (yhat)
      node[below right=-1pt, font=\footnotesize\bfseries] {$\hat{\bm{y}}$};

    % Decomposition parallelogram
    \coordinate (comp1end) at (1.35,0.26);
    \draw[colspace!50, thin, dashed] (comp1end) -- (yhat);
    \draw[colspace!50, thin, dashed] (0.45,1.34) -- (yhat);

    % Angle arc
    \draw[secondary, thin] (0.7,0.135) arc[start angle=11,end angle=72,radius=0.71];
    \node[secondary, font=\scriptsize] at (0.45,0.65) {$60°$};

    % Stable label
    \node[projection!70!black, font=\scriptsize, align=center] at (0,-0.8)
      {Unique, stable\\decomposition};
  \end{scope}

  % ============== RIGHT: Ill-conditioned ==============
  \begin{scope}[shift={(4.2,0)}]
    \node[secondary, font=\small\bfseries] at (0,3.6) {Ill-conditioned};

    % Column vectors nearly parallel
    \draw[->, colspace, very thick] (0,0) -- (2.5,1.0)
      node[right, font=\footnotesize\bfseries] {$\bm{x}_1$};
    \draw[->, colspace!70!black, very thick] (0,0) -- (2.3,1.2)
      node[above right, font=\footnotesize\bfseries] {$\bm{x}_2$};

    % Thin sliver region
    \fill[colspace!6] (0,0) -- (2.5,1.0) -- (2.3,1.2) -- cycle;

    % Projection yhat (stable)
    \coordinate (yhat) at (1.9,0.9);
    \draw[->, projection, thick] (0,0) -- (yhat)
      node[above left=-2pt, font=\footnotesize\bfseries] {$\hat{\bm{y}}$};

    % Decomposition 1
    \draw[response!60, thin, dashed] (0,0) -- (3.2,1.28);
    \draw[response!60, thin, dashed] (0,0) -- (-1.15,-0.6);
    \coordinate (d1a) at (3.2,1.28);
    \draw[response!50, thin, dashed] (d1a) -- (yhat);

    % Decomposition 2
    \draw[residual!50, thin, dashed] (0,0) -- (-0.5,-0.2);
    \draw[residual!50, thin, dashed] (0,0) -- (2.53,1.32);
    \coordinate (d2a) at (-0.5,-0.2);
    \draw[residual!50, thin, dashed] (d2a) -- (yhat);

    % Sliding annotation
    \draw[<->, secondary, thick] (2.8,1.15) -- (-0.8,-0.35);
    \node[secondary, font=\scriptsize, align=center, fill=white,
      inner sep=1pt] at (1.0,0.0)
      {coefficients\\slide};

    % Angle arc (small)
    \draw[secondary, thin] (0.6,0.24) arc[start angle=22,end angle=28,radius=0.64];
    \node[secondary, font=\scriptsize] at (1.0,0.55) {$\approx 5°$};

    % Unstable label
    \node[residual!70!black, font=\scriptsize, align=center] at (0,-0.8)
      {$\hat{\bm{y}}$ stable, but\\$\hat\beta_1,\hat\beta_2$ wildly unstable};
  \end{scope}

  % Dividing line
  \draw[secondary!30, thin] (0,-1.2) -- (0,3.8);

\end{tikzpicture}
\caption{Multicollinearity as geometry.
\textbf{Left}: well-separated predictors give a stable decomposition.
\textbf{Right}: nearly collinear predictors create a thin sliver; the
projection~$\hat{\bm{y}}$ is stable but the individual coefficients
slide dramatically along the sliver direction.}
\label{fig:multicollinearity}
\end{figure}


%======================================================================
\section{When the Geometry Fails}\label{sec:failures}
%======================================================================

\subsection{Multicollinearity: the geometry breaks}\label{sec:multicollinearity}

When predictors are nearly linearly dependent, $\Col(\bX)$ degenerates
from a ``thick'' plane into a paper-thin sliver.
Let us understand why this is dangerous and what exactly goes wrong.

\medskip
\noindent\textbf{What multicollinearity means.}
Multicollinearity occurs when one predictor can be approximately written as a linear combination of others.
For example, if $x_2 \approx 3x_1 + 2x_3$, then the three predictors contain nearly the same information.
The design matrix $\bX$ is close to being rank-deficient.

\medskip
\begin{remark}[Prerequisite: eigendecomposition of symmetric matrices]
Before proceeding, let us recall a fundamental fact from linear algebra.

Every real symmetric matrix $\bm{A} \in \R^{p \times p}$ can be
decomposed as
\[
  \bm{A} = \bQ \bm{\Lambda} \bQ^\top,
\]
where $\bQ = [\bm{q}_1 \mid \cdots \mid \bm{q}_p]$ is an orthogonal
matrix ($\bQ^\top \bQ = \bI$) whose columns are the \emph{eigenvectors},
and $\bm{\Lambda} = \diag(\lambda_1, \dots, \lambda_p)$ contains the
corresponding \emph{eigenvalues}.

Key properties:
\begin{itemize}[nosep]
  \item All eigenvalues of a real symmetric matrix are \emph{real}
        (not complex).
  \item The eigenvectors can be chosen to be \emph{orthonormal}.
  \item If $\bm{A}$ is also positive semi-definite (like $\bX^\top\bX$),
        all eigenvalues are non-negative: $\lambda_k \ge 0$.
  \item The eigenvalues measure the ``stretch'' of the matrix in each
        eigenvector direction: $\bm{A}\bm{q}_k = \lambda_k \bm{q}_k$.
\end{itemize}

The inverse (when it exists) is simply
$\bm{A}^{-1} = \bQ \bm{\Lambda}^{-1} \bQ^\top
= \bQ \, \diag(1/\lambda_1, \dots, 1/\lambda_p) \, \bQ^\top$.
This makes it clear that small eigenvalues produce large entries in the
inverse---the source of variance inflation under multicollinearity.
\end{remark}

\medskip
\noindent\textbf{Why it is dangerous: the pancake analogy.}
Consider fitting a model with two highly correlated predictors.
Geometrically, the column space $\Col(\bX)$ is a 2-dimensional plane in $\R^n$.
When the two predictors are nearly proportional, this plane is not a ``square'' plane---it is extremely elongated, like a thin pancake or a nearly 1-dimensional line.

Now imagine projecting the data vector $\by$ onto this thin pancake.
The projection itself is fine: $\yhat$ lands on the pancake and $\be$ is perpendicular.
The \emph{total} prediction $\yhat$ is stable.
But the \emph{individual coefficients} are not: there are many different ways to decompose $\yhat$ into contributions along two nearly parallel directions.
A tiny change in $\by$ can slide $\yhat$ along the pancake, dramatically changing the decomposition $(\hat{\beta}_1, \hat{\beta}_2)$ even though $\yhat = \hat{\beta}_1\bm{x}_1 + \hat{\beta}_2\bm{x}_2$ stays almost the same.

\medskip
\noindent\textbf{The eigenvalue perspective.}
The eigenvalue decomposition of $\bX^\top\bX$ makes this precise:
\[
  \bX^\top\bX = \bQ\bm{\Lambda}\bQ^\top,
  \qquad
  \bm{\Lambda} = \diag(\lambda_1,\dots,\lambda_p),
  \quad \lambda_1\ge\cdots\ge\lambda_p\ge 0.
\]
The eigenvalues measure the ``width'' of the column space in each principal direction.
A large eigenvalue $\lambda_k$ means the column space extends far in direction $\bm{q}_k$---the projection is stable there.
A small eigenvalue $\lambda_k$ means the column space is paper-thin in that direction---the projection's location is poorly determined.

The variance of the OLS coefficients is $\text{Var}(\bhat) = \sigma^2(\bX^\top\bX)^{-1} = \sigma^2\bQ\bm{\Lambda}^{-1}\bQ^\top$.
In the direction of the smallest eigenvalue, the variance is $\sigma^2/\lambda_p$---which explodes when $\lambda_p$ is small.

The \textbf{condition number} $\kappa = \lambda_1/\lambda_p$
measures the aspect ratio of the eigenvalue ellipsoid.
When $\kappa \gg 1$, the column space is elongated: the projection
onto the column space is stable (overall $R^2$ is fine), but the
projection's \emph{location within} the thin direction is wildly
unstable. Small perturbations in $\by$ cause large swings
in individual coefficients.

\medskip
\noindent\textbf{The key danger, summarised.}
Multicollinearity does \emph{not} make the overall model bad---$R^2$ and the predicted values $\yhat$ may be perfectly fine.
What it destroys is the interpretability of individual coefficients.
You cannot say ``a one-unit increase in $x_2$ leads to a $\hat{\beta}_2$ change in $y$, holding $x_1$ constant'' when $x_1$ and $x_2$ move together so tightly that ``holding $x_1$ constant'' is practically impossible.

\medskip
\noindent\textbf{Diagnostic:} The condition number $\kappa$ and the
eigenvalue spectrum. If $\kappa > 30$, investigate; if $\kappa > 100$,
individual coefficients may be unreliable.
Another common diagnostic is the \textbf{Variance Inflation Factor}
(VIF): $\text{VIF}_j = 1/(1-R_j^2)$, where $R_j^2$ is the $R^2$ from
regressing predictor $x_j$ on all other predictors.
$\text{VIF}_j > 10$ is often flagged as concerning.

\begin{misconception}
\textbf{``Multicollinearity makes the model bad.''}

Multicollinearity inflates standard errors of individual coefficients,
but it does \emph{not}:
\begin{itemize}[nosep]
  \item Bias the OLS estimates (they remain unbiased).
  \item Ruin the overall prediction ($\yhat$ and $R^2$ can be fine).
  \item Invalidate the $F$-test for the joint significance of a group
        of collinear predictors.
\end{itemize}
Multicollinearity is only a problem \emph{if you need to interpret
individual coefficients}.
If your goal is prediction, collinear predictors may cause no harm
at all---the instability in individual $\hat{\beta}_j$ averages out in
$\yhat = \bX\bhat$.
\end{misconception}

\subsection{Omitted variable bias: the geometry can't warn you}

\begin{warning}
Multicollinearity is a problem the geometry can \emph{see}.
Omitted variable bias is a problem the geometry \emph{cannot} see.
\end{warning}

Suppose the true model is
$\by = \bX_1\bm{\beta}_1 + \bX_2\bm{\beta}_2 + \bm{\varepsilon}$
but we estimate the ``short'' regression omitting $\bX_2$:
\[
  \bhat_1^{\text{short}}
  = (\bX_1^\top\bX_1)^{-1}\bX_1^\top\by
  = \bm{\beta}_1 + \underbrace{(\bX_1^\top\bX_1)^{-1}\bX_1^\top\bX_2}_{\bm{\Gamma}}\bm{\beta}_2
    + (\bX_1^\top\bX_1)^{-1}\bX_1^\top\bm{\varepsilon}.
\]
Taking expectations:
\begin{equation}\label{eq:ovb}
  \boxed{
  E[\bhat_1^{\text{short}}] = \bm{\beta}_1 + \bm{\Gamma}\bm{\beta}_2.
  }
\end{equation}
The bias $\bm{\Gamma}\bm{\beta}_2$ is the gap between projections onto
two different column spaces.
It exists whenever (i) $\bX_2$ is correlated with $\bX_1$
($\bm{\Gamma}\ne\bm{0}$) \emph{and} (ii) $\bX_2$ affects $\by$
($\bm{\beta}_2\ne\bm{0}$).

Both the short and long regressions have perfectly orthogonal residuals,
healthy condition numbers, and green diagnostics.
The geometry worked flawlessly for both projections.
The problem is not how you project---it is \emph{which wall} you
project onto.
Choosing the right column space requires \emph{causal reasoning},
which lies beyond the geometric framework.

\begin{remark}
Residual diagnostics cannot detect omitted variable bias. Both the
``short'' and ``long'' regressions produce residuals orthogonal to
their respective column spaces---by construction, not by virtue of
model correctness. The only safeguards against OVB are domain
knowledge and causal reasoning.
\end{remark}

\medskip
\noindent\textit{Multicollinearity makes the column space geometrically
thin, inflating coefficient variance.
Omitted variable bias makes the column space the wrong target entirely.
The first problem has a geometric remedy: can we \emph{constrain}
the projection to trade a little bias for a lot of stability?
This is the idea behind regularisation.}


%======================================================================
\section{Constraining the Shadow: Regularisation}\label{sec:regularisation}
%======================================================================

Until now, every result has been a consequence of orthogonal projection.
The hat matrix $\bH$ is idempotent and symmetric; the fitted values
$\yhat$ are the closest point in $\Col(\bX)$ to $\by$; the residual
$\be$ meets the column space at a right angle.
Regularisation breaks this pattern---Ridge and LASSO are \emph{not}
projections.
They deliberately sacrifice the ``closest point'' property to gain
stability.
This section explains how and why.

\begin{warning}
\textbf{Where the projection framework reaches its limits.}
The Ridge hat matrix
$\bH_{\text{ridge}} = \bX(\bX^\top\bX+\lambda\bI)^{-1}\bX^\top$
is \emph{not} idempotent
($\bH_{\text{ridge}}^2 \neq \bH_{\text{ridge}}$) and is \emph{not}
an orthogonal projection onto any subspace.
It does not ``drop a perpendicular''---it \emph{pulls the shadow
inward}, toward the origin, away from the foot of the perpendicular.

The LASSO solution is not even linear in $\by$.

The geometric language of shadows and right angles served us well for
OLS\@.  For regularisation, we need a different picture: that of
\textbf{constrained optimisation}, where the solution is determined
not by orthogonality but by the tangency between an objective function's
level sets and a constraint region.
\end{warning}

\medskip
\noindent\textbf{The telescope analogy.}
Think of the OLS estimate as looking through a telescope on a windy night.
The telescope points in the right direction on average (unbiasedness),
but the wind makes it shake violently (high variance from
multicollinearity or small sample size).
Each snapshot is centred on the truth, but individual frames are blurry
and useless.

Regularisation is like tightening the telescope's mount.
A tighter mount introduces a small systematic error---the telescope
now points slightly away from the true target (bias).
But the shaking is dramatically reduced (lower variance), so each
individual frame is much sharper and more useful.


%----------------------------------------------------------------------
\subsection{Ridge regression}
%----------------------------------------------------------------------

\begin{definition}
The \textbf{Ridge estimator} with penalty $\lambda\ge 0$ is
\[
  \bhat_{\text{ridge}}
  = \argmin_{\bm{\beta}}\left\{
      \norm{\by-\bX\bm{\beta}}^2 + \lambda\norm{\bm{\beta}}^2
    \right\}
  = (\bX^\top\bX + \lambda\bI)^{-1}\bX^\top\by.
\]
\end{definition}

\noindent\textbf{Eigenvalue shrinkage---the key mechanism.}
Ridge adds $\lambda\bI$ to $\bX^\top\bX$, which shifts every eigenvalue
$\lambda_k\mapsto\lambda_k+\lambda$.
Write the eigendecomposition $\bX^\top\bX = \bQ\bD\bQ^\top$, where
$\bD = \diag(\lambda_1,\dots,\lambda_p)$.
Then
\begin{equation}\label{eq:ridge-eigen}
  \bhat_{\text{ridge}}
  = \bQ\,\diag\!\left(\frac{\lambda_k}{\lambda_k+\lambda}\right)\bQ^\top
    \bhat_{\text{OLS}}.
\end{equation}
Each eigenvector component of $\bhat_{\text{OLS}}$ is multiplied by the
shrinkage factor $\lambda_k/(\lambda_k+\lambda)$, which lies in $[0,1]$.
\begin{itemize}[nosep]
  \item Directions with \emph{large} $\lambda_k$ (stable, well-determined
        by the data) are barely affected: the factor is close to $1$.
  \item Directions with \emph{small} $\lambda_k$ (unstable, dominated by
        noise) are aggressively shrunk toward zero: the factor is close
        to $0$.
\end{itemize}

\begin{keyinsight}
Ridge performs \emph{selective geometric surgery}: it identifies the
unstable directions of the column space via the eigenvalue spectrum
and shrinks those directions most aggressively, while leaving the
stable directions nearly intact.
This is \emph{anisotropic} shrinkage---it does not pull equally in
all directions, but preferentially collapses the dimensions where OLS
is least trustworthy.
\end{keyinsight}


%--- Example 6: Ridge shrinkage direction by direction -----------------

\begin{example}[Ridge Shrinkage Direction by Direction]\label{ex:ridge}
We work with a $2\times 2$ cross-product matrix $\bX^\top\bX$ whose
eigenvalues and eigenvectors are known, to see exactly how Ridge shrinks
each direction.

\medskip\noindent\textbf{Setup.}
Suppose (after centring) $\bX^\top\bX$ has eigendecomposition
$\bX^\top\bX = \bQ\bD\bQ^\top$ with eigenvalues $9$ and $1$:
\[
\bQ = \frac{1}{\sqrt{2}}\begin{pmatrix}1&1\\1&-1\end{pmatrix},\quad
\bD = \begin{pmatrix}9&0\\0&1\end{pmatrix}.
\]
This gives
\[
\bX^\top\bX
= \frac{1}{2}\begin{pmatrix}1&1\\1&-1\end{pmatrix}
\begin{pmatrix}9&0\\0&1\end{pmatrix}
\begin{pmatrix}1&1\\1&-1\end{pmatrix}
= \begin{pmatrix}5&4\\4&5\end{pmatrix}.
\]
The ``wide'' eigendirection is
$\bm{v}_1 = (1,1)^\top/\sqrt{2}$ (eigenvalue $9$, lots of data spread)
and the ``narrow'' eigendirection is
$\bm{v}_2 = (1,-1)^\top/\sqrt{2}$ (eigenvalue $1$, little spread).

Choose OLS solution $\bhat_{\text{OLS}} = (3,\;1)^\top$.
Decompose in the eigenbasis using unnormalised directions
$\bm{u}_1 = (1,1)^\top$, $\bm{u}_2 = (1,-1)^\top$:
\[
(3,1) = \alpha_1(1,1) + \alpha_2(1,-1)
\implies \alpha_1 = 2,\;\alpha_2 = 1.
\]
So OLS has component $2$ along the stable direction and $1$ along
the unstable direction.

\medskip\noindent\textbf{Ridge coefficients.}
In the eigenbasis, each component is simply multiplied by a shrinkage
factor:
\[
\alpha_k^{\text{Ridge}}
= \frac{\lambda_k}{\lambda_k + \lambda}\;\alpha_k^{\text{OLS}}.
\]
No matrix inversion is needed once you work in the eigenbasis.

\begin{center}
\renewcommand{\arraystretch}{1.4}
\begin{tabular}{c|cc|cc|cc|c}
\toprule
& \multicolumn{2}{c|}{Shrinkage factor}
& \multicolumn{2}{c|}{Eigencomponents}
& \multicolumn{2}{c|}{Coefficients} & \\
$\lambda$
& $\frac{9}{9+\lambda}$ & $\frac{1}{1+\lambda}$
& $\alpha_1^{\text{R}}$ & $\alpha_2^{\text{R}}$
& $\hat\beta_1^{\text{R}}$ & $\hat\beta_2^{\text{R}}$
& $\norm{\bhat_{\text{R}}}$ \\
\midrule
$0$   & $1$ & $1$
      & $2$ & $1$
      & $3$ & $1$
      & $\sqrt{10}\approx 3.16$ \\[2pt]
$1$   & $\frac{9}{10}$ & $\frac{1}{2}$
      & $\frac{9}{5}$ & $\frac{1}{2}$
      & $\frac{23}{10}$ & $\frac{13}{10}$
      & $2.64$ \\[2pt]
$10$  & $\frac{9}{19}$ & $\frac{1}{11}$
      & $\frac{18}{19}$ & $\frac{1}{11}$
      & $1.04$ & $0.86$
      & $1.35$ \\[2pt]
$100$ & $\frac{9}{109}$ & $\frac{1}{101}$
      & $\frac{18}{109}$ & $\frac{1}{101}$
      & $0.18$ & $0.16$
      & $0.24$ \\
\bottomrule
\end{tabular}
\end{center}

\medskip\noindent\textbf{The key pattern.}
At $\lambda=1$, the stable direction (eigenvalue $9$) retains $90\%$ of
its OLS component, while the unstable direction (eigenvalue $1$) is
already halved.
At $\lambda=10$, the unstable direction is crushed to $9\%$ of its OLS
value, while the stable direction still retains $47\%$.
As $\lambda$ increases, the coefficient vector traces a \emph{curved}
path from $(3,1)$ toward the origin---not a straight line, because
the unstable eigencomponent collapses first.
\end{example}


\medskip
\noindent\textbf{Geometric interpretation.}
The Ridge solution is the point where the smallest RSS contour
ellipsoid is tangent to the $L^2$ sphere
$\|\bm{\beta}\|^2 \le t$ (Figure~\ref{fig:ridge-lasso}).
Because the sphere has no corners, the tangent point generically
has all coefficients nonzero.


%----------------------------------------------------------------------
\subsection{LASSO regression}
%----------------------------------------------------------------------

\begin{definition}
The \textbf{LASSO estimator} with penalty $\lambda\ge 0$ is
\[
  \bhat_{\text{lasso}}
  = \argmin_{\bm{\beta}}\left\{
      \norm{\by-\bX\bm{\beta}}^2 + \lambda\norm{\bm{\beta}}_1
    \right\},
\]
where $\norm{\bm{\beta}}_1 = \sum_j|\beta_j|$.
\end{definition}

\noindent\textbf{Corner geometry and sparsity.}
The $L^1$ constraint region
$\{(\beta_1,\dots,\beta_p) : |\beta_1|+\cdots+|\beta_p| \le t\}$
is a diamond (cross-polytope) instead of a sphere.
Its surface has \emph{corners}---points where one or more coordinates
are exactly zero.
The RSS contour ellipses are more likely to first touch the diamond
at a corner, producing \textbf{exact zeros} in the coefficient vector
(see Figure~\ref{fig:ridge-lasso}, dashed diamond).

Why?
A corner is a ``pointy'' protrusion: an expanding ellipse approaching
from the direction of $\bhat_{\text{OLS}}$ will generically hit a
corner before it hits a smooth edge.
In contrast, the $L^2$ sphere has no corners, so the tangency point
generically has \emph{all} coordinates nonzero.

\medskip
\noindent\textbf{Sparsity as a corner solution.}
LASSO performs automatic variable selection: as $\lambda$ increases,
coefficients are driven to zero one by one.
Each zero corresponds to the solution hitting a corner of the $L^1$
diamond.
This is not an algebraic trick---it is a consequence of the geometry
of touching a polytope with an ellipsoid.

\begin{stopandthink}
\textbf{Why corners attract LASSO solutions.}

Grab a piece of paper and try this:
\begin{enumerate}[nosep]
  \item In $\R^2$, draw the $L^1$ diamond
        $\{(\beta_1, \beta_2) : |\beta_1| + |\beta_2| \le 1\}$.
        It has corners at $(1,0)$, $(-1,0)$, $(0,1)$, $(0,-1)$.
  \item Draw a family of concentric ellipses centred at some point
        outside the diamond (this represents the RSS contours centred
        at $\bhat_{\text{OLS}}$).
  \item Starting from the smallest ellipse, gradually expand it until
        it first touches the diamond. Where does it touch?
  \item Now replace the diamond with the $L^2$ circle
        $\{(\beta_1, \beta_2) : \beta_1^2 + \beta_2^2 \le 1\}$
        and repeat. Where does the first contact occur now?
\end{enumerate}
\emph{Key observation:}
The diamond's corners are ``pointy''---an expanding ellipse is far more
likely to first touch a corner (where one coordinate is exactly zero)
than a smooth edge.
The circle has no corners, so the first contact is almost always at a
point where \emph{both} coordinates are nonzero.
This is why LASSO produces sparse solutions (exact zeros) and Ridge
does not.
The geometry of the constraint set, not the algebra, is what drives
sparsity.
\end{stopandthink}


%----------------------------------------------------------------------
\subsection{Ridge vs.\ LASSO: a comparison}
%----------------------------------------------------------------------

Figure~\ref{fig:ridge-lasso} illustrates both constraint regions
simultaneously.
The OLS estimate $\bhat_{\text{OLS}}$ sits at the centre of the
elliptical RSS contours.
Ridge constrains the solution to the $L^2$ sphere (solid circle);
LASSO constrains it to the $L^1$ diamond (dashed).
The regularised solutions are the tangency points.

\begin{center}
\begin{tabular}{lcc}
\toprule
& \textbf{Ridge ($L^2$)} & \textbf{LASSO ($L^1$)} \\
\midrule
Constraint shape & Sphere (smooth) & Diamond (corners) \\
Shrinkage type & Anisotropic via eigenvalues & Along coordinate axes \\
Coefficient path & Smooth shrinkage toward 0 & Coefficients hit 0 exactly \\
Variable selection & No (all coefficients remain) & Yes (sparsity) \\
Best when & Many small effects (dense) & Few large effects (sparse) \\
Closed-form solution & Yes & No (requires iterative solver) \\
\bottomrule
\end{tabular}
\end{center}

% --- Figure 5: Ridge vs LASSO constraint regions ---

\begin{figure}[ht]
\centering
\begin{tikzpicture}[>=Stealth, scale=1.15]

  % --- Axes ---
  \draw[->, secondary, thick] (-3.2,0) -- (4.2,0) node[right, font=\small] {$\beta_1$};
  \draw[->, secondary, thick] (0,-3.2) -- (0,4.2) node[above, font=\small] {$\beta_2$};

  % --- OLS solution ---
  \coordinate (ols) at (2.8,2.4);
  \fill[response] (ols) circle (3pt);
  \node[above right, response!80!black, font=\footnotesize\bfseries] at (ols)
    {$\hat{\bm{\beta}}_{\mathrm{OLS}}$};

  % --- RSS contour ellipses (centered on OLS, tilted) ---
  \foreach \r in {0.6, 1.2, 1.85, 2.6, 3.4} {
    \draw[secondary!35, thin, rotate around={25:(ols)}]
      (ols) ellipse ({\r} and {\r*0.45});
  }
  \node[secondary, font=\scriptsize] at (3.8,0.6) {RSS contours};

  % --- L2 constraint region (circle) ---
  \draw[constraint, thick, fill=constraint!8] (0,0) circle (1.5);
  \node[constraint!80!black, font=\scriptsize] at (-1.0,-1.8)
    {$\beta_1^2+\beta_2^2 \le t$};

  % --- L1 constraint region (diamond) ---
  \draw[constraint, thick, dashed, fill=constraint!4]
    (1.9,0) -- (0,1.9) -- (-1.9,0) -- (0,-1.9) -- cycle;
  \node[constraint!80!black, font=\scriptsize] at (1.7,-1.8)
    {$|\beta_1|+|\beta_2| \le t$};

  % --- Ridge solution (ellipse tangent to circle) ---
  \coordinate (ridge) at ({1.5*cos(40)},{1.5*sin(40)});
  \fill[projection] (ridge) circle (3pt);
  \node[above left, projection!80!black, font=\footnotesize\bfseries] at (ridge)
    {$\hat{\bm{\beta}}_{\mathrm{Ridge}}$};
  % Ridge ellipse (approximate tangent)
  \draw[projection!60, thin, rotate around={25:(ols)}]
    (ols) ellipse ({2.05} and {2.05*0.45});

  % --- LASSO solution (ellipse tangent to diamond corner) ---
  \coordinate (lasso) at (1.9,0);
  \fill[residual] (lasso) circle (3pt);
  \node[below right, residual!80!black, font=\footnotesize\bfseries] at (lasso)
    {$\hat{\bm{\beta}}_{\mathrm{LASSO}}$};
  % LASSO ellipse (approximate tangent at corner)
  \draw[residual!50, thin, rotate around={25:(ols)}]
    (ols) ellipse ({1.65} and {1.65*0.45});

  % --- Annotation: why corners produce zeros ---
  \draw[->, secondary, thin] (3.0,-2.6) -- ($(lasso)+(0.15,-0.15)$);
  \node[secondary, font=\scriptsize, align=left, text width=3.2cm] at (3.5,-3.2)
    {Diamond corner lies on axis:\\$\hat\beta_2 = 0$ (sparsity!)};

  % --- Annotation: circle has no corners ---
  \draw[->, secondary, thin] (-2.8,2.2) -- ($(ridge)+(-0.15,0.1)$);
  \node[secondary, font=\scriptsize, align=left, text width=3cm] at (-3.2,2.8)
    {Circle is smooth: solution\\generically has $\hat\beta_j\neq 0$};

\end{tikzpicture}
\caption{Ridge vs.\ LASSO constraint regions.  The OLS
estimate~$\hat{\bm{\beta}}_{\mathrm{OLS}}$ (amber dot) sits at the
centre of the elliptical RSS contours.  Regularisation constrains the
solution to lie within a region around the origin.  \textbf{Ridge}
(solid circle, $L^2$) shrinks all coefficients smoothly toward zero.
\textbf{LASSO} (dashed diamond, $L^1$) has \emph{corners} on the
coordinate axes; the first contour ellipse to touch the diamond
typically hits a corner, setting some coefficients exactly to zero
and producing a sparse model.}
\label{fig:ridge-lasso}
\end{figure}


%----------------------------------------------------------------------
\subsection{Reconnecting to the geometric framework}
%----------------------------------------------------------------------

Neither Ridge nor LASSO is an orthogonal projection.
Both sacrifice the ``closest point'' property---accepting bias---in
exchange for a (potentially large) reduction in variance.
The telescope analogy captures the trade-off:
you give up perfect aim to eliminate the shake.

Yet the connection to the projection framework runs deeper than it
first appears.
Recall from Section~\ref{sec:multicollinearity} that
multicollinearity makes the column space geometrically thin:
the eigenvalue spectrum of $\bX^\top\bX$ has a few large values
and some dangerously small ones.
OLS faithfully projects onto this flattened space, and in doing so
amplifies noise along the thin directions.
Ridge responds by \emph{reshaping} the eigenvalue landscape:
it replaces $\lambda_k$ with $\lambda_k + \lambda$, inflating every
eigenvalue and rounding out the column space's geometry.
The thinner a direction was, the more it gets inflated in relative
terms---precisely the directions that needed stabilising.

Ridge is the projection framework's answer to its own weakness.
When the geometry is unstable, you reshape the geometry.


%--- Gaussian prerequisite ---

\medskip
\begin{remark}[Prerequisite: where the Gaussian enters]
Everything up to this point has been \emph{purely geometric}---no
distributional assumptions were needed.
The Projection Theorem, the hat matrix properties, Pythagoras, FWL,
leverage, influence, multicollinearity, OVB, and even regularisation
all follow from linear algebra alone.

To perform \emph{inference}---hypothesis tests, $p$-values, confidence
intervals---we need a probability model.
The standard assumption is:
\[
  \by = \bX\bm{\beta} + \bm{\varepsilon},
  \qquad
  \bm{\varepsilon} \sim \mathcal{N}(\bm{0},\, \sigma^2 \bI_n).
\]
That is, the errors are independent, identically distributed, and
Gaussian (normally distributed).

Why does the Gaussian matter?
\begin{itemize}[nosep]
  \item \textbf{Exactness:} Under Gaussianity, $\bhat$ is exactly
        $\mathcal{N}(\bm{\beta},\, \sigma^2(\bX^\top\bX)^{-1})$---not
        just approximately. This gives us exact $t$- and $F$-distributions
        for test statistics, valid in any sample size.
  \item \textbf{Without Gaussianity:} The Central Limit Theorem still
        ensures that $\bhat$ is \emph{approximately} normal in large samples.
        So inference is approximately valid even when the true error
        distribution is non-Gaussian---but only asymptotically.
  \item \textbf{Geometry of the Gaussian:} A Gaussian vector projected
        onto a subspace remains Gaussian. This is why the hat matrix
        preserves normality: $\yhat = \bH\by$ and
        $\be = \bM\by$ are both Gaussian, and they are
        \emph{independent} (because $\bH\bM = \bm{0}$).
        Independence of $\yhat$ and $\be$ is what makes the $F$-test work.
\end{itemize}
The next section uses these distributional facts to build hypothesis tests.
\end{remark}

\medskip
\noindent\textit{Regularisation reshapes the geometry to stabilise
estimates, but it cannot tell us whether the extra predictors are
\emph{statistically significant}.
For that, we need to compare two projections---one with and one
without the extra variables---and ask whether the improvement is
larger than what noise alone would produce.
This brings us to the $F$-test.}


%======================================================================
\section{Testing with Geometry: The $F$-Test}\label{sec:ftest}
%======================================================================

\subsection{Comparing two projections}

Consider two nested models:
\begin{itemize}[nosep]
  \item \textbf{Restricted:} $\by = \bX_R\bm{\beta}_R + \be_R$,
        with column space $\Col(\bX_R)$.
  \item \textbf{Unrestricted:} $\by = \bX\bm{\beta} + \be$,
        with column space $\Col(\bX)\supseteq\Col(\bX_R)$.
\end{itemize}
Both models project $\by$ onto their respective column spaces.
The restricted projection $\yhat_R$ and the unrestricted projection
$\yhat$ satisfy $\norm{\be_R}\ge\norm{\be}$ (by Proposition~\ref{prop:r2mono}).
The gap $\yhat - \yhat_R$ lies in $\Col(\bX)$ and is perpendicular to
$\Col(\bX_R)$.

\subsection{The $F$-statistic}

\begin{definition}
The $F$-statistic for testing $H_0$: the $q$ extra predictors
have zero coefficients is
\begin{equation}\label{eq:fstat}
  F = \frac{(\text{SSE}_R - \text{SSE})/q}{\text{SSE}/(n-p)}
    = \frac{\norm{\yhat-\yhat_R}^2/q}{\norm{\be}^2/(n-p)}.
\end{equation}
Under $H_0$ and standard Gaussian assumptions,
$F\sim F_{q,\,n-p}$.
\end{definition}

\begin{keyinsight}
The $F$-test measures whether the \emph{distance between two projections}
is surprisingly large relative to the residual noise.
The numerator is the squared length of the gap between projections,
normalised by the number of extra dimensions ($q$).
The denominator is the residual noise level.
If $F\approx 1$, the extra predictors captured only noise.
If $F\gg 1$, they captured real signal.
\end{keyinsight}

% --- Figure 6: F-test gap between projections ---

\begin{figure}[ht]
\centering
\tdplotsetmaincoords{62}{125}
\begin{tikzpicture}[tdplot_main_coords, scale=1.5, >=Stealth]

  % --- Full plane Col(X) ---
  \fill[colspace!10] (-3,-2.5,0) -- (3.2,-2.5,0) -- (3.2,2.5,0) -- (-3,2.5,0) -- cycle;
  \draw[colspace!30, line width=0.4pt]
    (-3,-2.5,0) -- (3.2,-2.5,0) -- (3.2,2.5,0) -- (-3,2.5,0) -- cycle;
  \node[colspace!70!black, font=\scriptsize\itshape] at (2.7,-2.0,0)
    {$\operatorname{Col}(\bm{X})$};

  % --- Restricted subspace Col(X_R) as a line in the plane ---
  \draw[constraint!70, very thick] (-2.8,-0.56,0) -- (2.8,0.56,0);
  \node[constraint!80!black, font=\scriptsize\itshape] at (2.9,0.9,0)
    {$\operatorname{Col}(\bm{X}_R)$};

  % Origin
  \fill[secondary] (0,0,0) circle (1.2pt);
  \node[below left, secondary, font=\scriptsize] at (0,0,0) {$\bm{0}$};

  % --- y above the plane ---
  \coordinate (y) at (1.6,1.4,2.6);
  \draw[->, response, very thick] (0,0,0) -- (y)
    node[above right, font=\footnotesize\bfseries] {$\bm{y}$};

  % --- yhat: projection onto full Col(X) ---
  \coordinate (yhat) at (1.6,1.4,0);
  \draw[->, projection, very thick] (0,0,0) -- (yhat)
    node[below right=1pt, font=\footnotesize\bfseries] {$\hat{\bm{y}}$};

  % --- yhatR: projection onto Col(X_R) (on the line) ---
  \coordinate (yhatR) at (1.75,0.35,0);
  \draw[->, constraint, thick] (0,0,0) -- (yhatR)
    node[below=3pt, font=\footnotesize\bfseries] {$\hat{\bm{y}}_R$};

  % --- Residual e = y - yhat (perpendicular to plane, going up) ---
  \draw[->, residual, thick, dashed] (yhat) -- (y);
  \node[right=3pt, residual!80!black, font=\footnotesize] at ($(yhat)!0.5!(y)$)
    {$\bm{e}=\bm{y}-\hat{\bm{y}}$};

  % --- Gap: yhat - yhatR (in the plane, perp to Col(X_R)) ---
  \draw[->, response!80!black, thick] (yhatR) -- (yhat);
  \node[above=3pt, response!80!black, font=\footnotesize] at ($(yhatR)!0.5!(yhat)$)
    {$\hat{\bm{y}}-\hat{\bm{y}}_R$};

  % --- Right angle 1: e perpendicular to plane at yhat ---
  \draw[secondary, line width=0.5pt]
    ($(yhat)+(-0.16,0,0)$) -- ($(yhat)+(-0.16,0,0.16)$) -- ($(yhat)+(0,0,0.16)$);

  % --- Right angle 2: gap perpendicular to Col(X_R) at yhatR ---
  \pgfmathsetmacro{\ux}{0.98}
  \pgfmathsetmacro{\uy}{0.196}
  \pgfmathsetmacro{\vx}{-0.196}
  \pgfmathsetmacro{\vy}{0.98}
  \draw[secondary, line width=0.5pt]
    ($(yhatR)+0.14*(\ux,\uy,0)$)
    -- ($(yhatR)+0.14*(\ux,\uy,0)+0.14*(\vx,\vy,0)$)
    -- ($(yhatR)+0.14*(\vx,\vy,0)$);

  % --- Right angle 3: (yhat - yhatR) perpendicular to e ---
  \draw[secondary, line width=0.5pt]
    ($(yhat)+(0,-0.16,0)$) -- ($(yhat)+(0,-0.16,0.16)$) -- ($(yhat)+(0,0,0.16)$);

  % --- F-statistic annotation ---
  \draw[decorate, decoration={brace, amplitude=4pt, mirror},
    secondary, thick]
    ($(yhatR)+(0.2,-0.15,0)$) -- ($(yhat)+(0.2,-0.15,0)$);
  \node[secondary, font=\scriptsize, right, align=left] at (2.4,0.2,0)
    {$\|\hat{\bm{y}}-\hat{\bm{y}}_R\|^2/(p-q)$\\[2pt]
     $F$-numerator};

  \draw[decorate, decoration={brace, amplitude=4pt},
    secondary, thick]
    ($(yhat)+(0.35,0,0)$) -- ($(y)+(0.35,0,0)$);
  \node[secondary, font=\scriptsize, right, align=left] at (2.3,0.7,1.3)
    {$\|\bm{e}\|^2/(n-p)$\\[2pt]
     $F$-denominator};

\end{tikzpicture}
\caption{The $F$-test as the gap between two projections.
$\hat{\bm{y}}_R$ is the projection of~$\bm{y}$ onto the restricted
subspace $\operatorname{Col}(\bm{X}_R)$ (a line), and $\hat{\bm{y}}$ is
the projection onto the full $\operatorname{Col}(\bm{X})$ (a plane
containing the line).  The three segments are mutually perpendicular by
the Pythagorean theorem:
$\|\bm{y}-\hat{\bm{y}}_R\|^2 =
 \|\hat{\bm{y}}-\hat{\bm{y}}_R\|^2 + \|\bm{e}\|^2$.
The $F$-statistic compares the ``explained gap''
$\|\hat{\bm{y}}-\hat{\bm{y}}_R\|^2$ (what the extra predictors bought)
to the residual $\|\bm{e}\|^2$ (what remains unexplained), each
normalised by its degrees of freedom.}
\label{fig:f-test}
\end{figure}

\begin{remark}
When $q=1$ (testing a single coefficient), $F = t^2$ where $t$ is the
usual $t$-statistic.
The $t$-test is the $F$-test in one dimension.
\end{remark}

\subsection{Confidence ellipsoids}

The sampling distribution of $\bhat$ is
$\bhat\sim\mathcal{N}(\bm{\beta},\;\sigma^2(\bX^\top\bX)^{-1})$.
A $(1-\alpha)$-confidence ellipsoid for $\bm{\beta}$ is
\[
  \bigl\{\bm{b}\in\R^p :
  (\bhat-\bm{b})^\top(\bX^\top\bX)(\bhat-\bm{b})
  \le p\,s^2\,F_{p,n-p,\alpha}
  \bigr\}.
\]
The ellipsoid's axes align with the eigenvectors of $\bX^\top\bX$
and its semi-axis lengths are proportional to $1/\sqrt{\lambda_k}$.
Directions where the column space is ``thin'' (small $\lambda_k$)
produce long axes---more uncertainty.
This is the same eigenvalue structure that drives multicollinearity
and Ridge shrinkage: everything connects.

\medskip
\noindent\textit{The $F$-test and confidence ellipsoids complete our
inferential toolkit.
But before we summarise, it is worth stepping back and asking:
where does the geometric framework \emph{break down}?
What can projections not tell us, and where must we look beyond
linear algebra for answers?}


%======================================================================
\section{The Limits of the Geometric Framework}\label{sec:limits}
%======================================================================

The Projection Theorem is powerful but not omniscient.
Three fundamental limitations bound the geometric framework, and
intellectual honesty demands that we name them clearly.

\subsection{Causality: the geometry cannot choose the wall}

The geometry tells you \emph{how well} you are projecting.
It cannot tell you \emph{whether you are projecting onto the right wall}.
Omitted variable bias (\S\ref{sec:failures}) is invisible to every
geometric diagnostic: the residuals are still perpendicular, $R^2$ still
decomposes via Pythagoras, the hat matrix is still idempotent.
The Scoreboard reads green for a perfectly projected regression on the
wrong column space.

To choose the right wall, you need \emph{causal reasoning}: directed
acyclic graphs, potential outcomes, natural experiments, instrumental
variables. These are structural tools, not geometric ones.
No amount of staring at shadows will tell you whether the flashlight
is aimed at the right surface.

\subsection{The inner product: whose notion of ``closest''?}

Everything in these notes assumed the Euclidean inner product---all
observations weighted equally, all errors independent.
If observations have \textbf{heteroscedastic} errors
($\text{Var}(\varepsilon_i) = \sigma_i^2$) or are
\textbf{autocorrelated},
the inner product is wrong.
The ``perpendicular'' is not truly perpendicular in the correct metric,
and the ``closest point'' is not truly closest.

\textbf{Generalised Least Squares (GLS)} fixes this by replacing the
Euclidean inner product $\ip{\bm{u}}{\bm{v}}=\bm{u}^\top\bm{v}$ with
$\ip{\bm{u}}{\bm{v}}_\Omega = \bm{u}^\top\bm{\Omega}^{-1}\bm{v}$,
where $\bm{\Omega} = \text{Var}(\bm{\varepsilon})$.
The Projection Theorem still applies, but in a different geometry---a
tilted space where ``perpendicular'' has a new meaning:
\[
  \bhat_{\text{GLS}}
  = (\bX^\top\bm{\Omega}^{-1}\bX)^{-1}\bX^\top\bm{\Omega}^{-1}\by.
\]
The moral: the geometry generalises, but only if you choose the right
inner product.

\subsection{Linearity: the column space is flat}

The column space is \emph{flat}---a hyperplane through the origin.
If the true relationship is curved, the best flat projection may be
a poor representation, like pressing a globe onto a table.
Polynomial features enlarge the column space to include some curvature,
but the bias-variance trade-off (Notebook~3's overfitting trap) limits
how far this can go.
Kernel methods, trees, and neural networks define curved ``column
spaces''---function spaces that lie beyond the reach of linear
projection geometry.

The Projection Theorem is not wrong in these settings.  It is simply
not enough.


%======================================================================
\section{Summary: The Geometric Rosetta Stone}\label{sec:rosetta}
%======================================================================

\subsection*{Closing the circle}

In 1801, Carl Friedrich Gauss took 41 noisy astronomical observations and
dropped a perpendicular from them onto the subspace of Keplerian orbits.
The shadow---the nearest point on that subspace---told him where to
point the telescope.  On New Year's Eve, von~Zach found Ceres within
half a degree of the predicted position.

Every time you call \texttt{lm()} in R or
\texttt{sklearn.LinearRegression().fit()} in Python, you perform
the same act that Gauss performed: you cast a shadow of noisy data
onto a subspace of possible fitted values, you measure the
perpendicular gap that remains, and you trust the right angle.

The mathematics has not changed in two centuries.  Only the notation
is different.

\subsection*{The Rosetta Stone}

Every concept in OLS regression has a geometric face and an algebraic
face.  Table~\ref{tab:rosetta} provides the complete correspondence.
Where the Pythagorean decomposition and $R^2$ require demeaned vectors,
this is noted explicitly.

\begin{table}[ht]
\centering
\caption{The Geometric Rosetta Stone.}\label{tab:rosetta}
\small
\begin{tabular}{@{}lll@{}}
\toprule
\textbf{Statistical concept} & \textbf{Geometric interpretation} & \textbf{Formula / characterisation} \\
\midrule
\multicolumn{3}{@{}l}{\textit{Core projection}} \\[2pt]
Fitted values $\yhat$ & Shadow on the wall & $\bH\by$ \\
Residuals $\be$ & Perpendicular gap & $(\bI-\bH)\by$ \\
Coefficients $\bhat$ & Coordinates of the shadow & $(\bX^\top\bX)^{-1}\bX^\top\by$ \\
Normal equations & Orthogonality condition & $\bX^\top\be=\bm{0}$ \\
Intercept $\hat{\beta}_0$ & Ensures $\bar{e}=0$; centres the & $\bm{1}\in\Col(\bX)$ guarantees \\
 & projection around the data mean & $\be\perp\bm{1}$ \\[4pt]
\midrule
\multicolumn{3}{@{}l}{\textit{Decomposition and fit quality (demeaned vectors $\tilde{\by}=\by-\bar{y}\bm{1}$, $\tilde{\yhat}=\yhat-\bar{y}\bm{1}$)}} \\[2pt]
SST $=$ SSR $+$ SSE & Pythagorean theorem & $\norm{\tilde{\by}}^2=\norm{\tilde{\yhat}}^2+\norm{\be}^2$ \\
$R^2$ & Squared cosine of angle & $\cos^2\theta = \norm{\tilde{\yhat}}^2/\norm{\tilde{\by}}^2$ \\
Degrees of freedom & Dimension of the residual & $\dim\bigl(\Col(\bX)^\perp\bigr) = n-p$ \\
 & subspace & \\
Adjusted $R^2$ & $R^2$ corrected for dimension & $1 - \frac{\norm{\be}^2/(n-p)}{\norm{\tilde{\by}}^2/(n-1)}$ \\[4pt]
\midrule
\multicolumn{3}{@{}l}{\textit{Projection operators and diagnostics}} \\[2pt]
Hat matrix $\bH$ & Projection operator & $\bH^2=\bH,\;\bH^\top=\bH$ \\
Leverage $h_{ii}$ & Arm length in column space & $\diag(\bH)$ \\
Cook's distance $D_i$ & Arm length $\times$ surprise & $\propto h_{ii}\cdot e_i^2$ \\
Gauss--Markov (BLUE) & Closest point in subspace & Projection Theorem \\[4pt]
\midrule
\multicolumn{3}{@{}l}{\textit{Structure and partial effects}} \\[2pt]
FWL theorem & Sequential orthogonal projections & $\bhat_2 = (\bX_2^\top\bM_1\bX_2)^{-1}\bX_2^\top\bM_1\by$ \\
Multicollinearity & Thin, nearly degenerate column space & $\kappa(\bX^\top\bX)\gg 1$ \\
OVB & Wrong column space & $\text{Bias} = \bm{\Gamma}\bm{\beta}_2$ \\[4pt]
\midrule
\multicolumn{3}{@{}l}{\textit{Regularisation (constrained optimisation: trading bias for stability)}} \\[2pt]
Ridge ($L^2$) & Constrained optimisation within a sphere; & Shrinks unstable eigenvector \\
 & trades bias for variance reduction & directions: $\lambda_k/(\lambda_k+\lambda)$ \\
LASSO ($L^1$) & Constrained optimisation within a diamond; & Corner solutions $\Rightarrow$ exact zeros \\
 & trades bias for sparsity & $\Rightarrow$ automatic variable selection \\[4pt]
\midrule
\multicolumn{3}{@{}l}{\textit{Inference}} \\[2pt]
$F$-test & Gap between two projections & $F = \frac{\norm{\yhat-\yhat_R}^2/q}{\norm{\be}^2/(n-p)}$ \\
Confidence ellipsoid & Region of plausible shadows & Axes $\propto 1/\sqrt{\lambda_k}$ \\
\bottomrule
\end{tabular}
\end{table}

\subsection*{The story in sixteen lines}

\begin{center}
\fbox{\parbox{0.88\textwidth}{\centering\itshape
Regression is a projection.\\
Coefficients are coordinates of the shadow.\\
The residual is perpendicular to the column space.\\
The intercept guarantees the residuals have mean zero.\\
$R^2$ is the squared cosine of an angle between demeaned vectors.\\
Variance decomposition is Pythagoras applied to a right triangle.\\
Degrees of freedom count the dimensions the residual can roam.\\
Gauss--Markov says: no linear unbiased estimator gets closer.\\
Controlling for a variable means projecting it away first.\\
Leverage measures how far each point reaches into the column space.\\
Influence is leverage times surprise.\\
Multicollinearity flattens the column space to a sliver.\\
Omitted variables tilt the column space to the wrong wall.\\
The $F$-test measures the gap between two nested projections.\\
Ridge and LASSO trade bias for stability---they are not projections.\\
And geometry alone cannot tell you which wall to aim at.
}}
\end{center}

\bigskip

\noindent
Gauss pointed a pencil at the sky and found an asteroid.  You point a
line of code at a dataset and find a pattern.  In both cases, the
method is the same: drop a perpendicular, read the shadow, respect the
right angle.

\medskip

\noindent
The geometry was always there, hiding inside every formula you
memorised and every $p$-value you reported.  These notes did not
create it.  They just turned on the light.

%======================================================================
% References
%======================================================================

\bigskip
\noindent\textbf{Further reading.}
\begin{itemize}[nosep]
  \item S.\,M.\ Stigler, \emph{The History of Statistics: The Measurement of
        Uncertainty before 1900}, Harvard University Press, 1986.
        \textit{Chapter~4 gives a superb account of Gauss, Ceres, and the
        birth of least squares.}
  \item G.\,W.\ Dunnington, \emph{Gauss: Titan of Science}, Mathematical
        Association of America, 2004 (reprint).
        \textit{The definitive biography; see chapters on the Ceres
        prediction and the development of the method of least squares.}
  \item G.\,A.\,F.\ Seber, \emph{Linear Regression Analysis}, Wiley, 1977.
  \item R.\ Davidson and J.\,G.\ MacKinnon, \emph{Estimation and Inference
        in Econometrics}, Oxford University Press, 1993.
  \item R.\ Christensen, \emph{Plane Answers to Complex Questions: The Theory
        of Linear Models}, 4th ed., Springer, 2011.
  \item S.\ Weisberg, \emph{Applied Linear Regression}, 4th ed., Wiley, 2014.
\end{itemize}


%======================================================================
\appendix
\section{Proofs for the Mathematically Curious}\label{app:proofs}
%======================================================================

These proofs are included for completeness and for readers who want
to verify the foundations. They are not required for following the
main narrative.


%----------------------------------------------------------------------
\subsection{Proof of the Projection Theorem (Theorem~\ref{thm:projection})}
%----------------------------------------------------------------------

\begin{proof}
We work in $V = \R^n$ with the Euclidean inner product, and $S$ is a
$p$-dimensional subspace (hence closed).  The proof has three parts:
existence, uniqueness, and the orthogonality characterisation.

\medskip
\noindent\textbf{Step 1: Existence.}
Let $d = \inf_{\bm{s}\in S}\norm{\by - \bm{s}}$.
Choose any $\bm{s}_0 \in S$. Then every minimising candidate satisfies
$\norm{\by - \bm{s}} \le \norm{\by - \bm{s}_0}$,
so the search is confined to the closed, bounded set
$\{\bm{s}\in S : \norm{\by - \bm{s}} \le \norm{\by - \bm{s}_0}\}$.
Since $S$ is a subspace of a finite-dimensional space, this set is compact.
The continuous function $\bm{s}\mapsto\norm{\by - \bm{s}}$ attains its
minimum on a compact set, so there exists $\yhat\in S$ with
$\norm{\by - \yhat} = d$.

\medskip
\noindent\textbf{Step 2: Uniqueness.}
Suppose $\yhat$ and $\yhat'$ are both minimisers.
By the parallelogram law,
\[
  \norm{(\by - \yhat) + (\by - \yhat')}^2
  + \norm{(\by - \yhat) - (\by - \yhat')}^2
  = 2\norm{\by - \yhat}^2 + 2\norm{\by - \yhat'}^2.
\]
The left-hand side rewrites as
$\norm{2\by - \yhat - \yhat'}^2 + \norm{\yhat' - \yhat}^2$,
and the right-hand side equals $4d^2$.
Now $\frac{1}{2}(\yhat + \yhat') \in S$ (since $S$ is a subspace), so
\[
  \norm{2\by - \yhat - \yhat'}^2
  = 4\,\norm{\by - \tfrac{\yhat+\yhat'}{2}}^2
  \ge 4d^2.
\]
Therefore $\norm{\yhat' - \yhat}^2 \le 4d^2 - 4d^2 = 0$,
which forces $\yhat' = \yhat$.

\medskip
\noindent\textbf{Step 3: Orthogonality condition.}
Let $\bm{v}\in S$ be arbitrary and $t\in\R$.
Since $\yhat + t\bm{v}\in S$, the minimality of $\yhat$ gives
\[
  f(t)
  \;=\; \norm{\by - \yhat - t\bm{v}}^2
  \;=\; \norm{\by - \yhat}^2
        - 2t\,\ip{\by - \yhat}{\bm{v}}
        + t^2\norm{\bm{v}}^2
  \;\ge\; f(0)
  \quad\text{for all } t\in\R.
\]
The function $f$ is a quadratic in $t$ with minimum at $t=0$.
Setting $f'(0) = 0$:
\[
  f'(0) = -2\,\ip{\by-\yhat}{\bm{v}} = 0
  \qquad\Longrightarrow\qquad
  \ip{\by - \yhat}{\bm{v}} = 0.
\]
Since $\bm{v}\in S$ was arbitrary, $(\by - \yhat)\perp S$.

\medskip
\noindent\textbf{Conversely}, if $\yhat\in S$ satisfies $(\by-\yhat)\perp S$,
then for any $\bm{s}\in S$:
\[
  \norm{\by - \bm{s}}^2
  = \norm{(\by - \yhat) + (\yhat - \bm{s})}^2
  = \norm{\by - \yhat}^2 + \norm{\yhat - \bm{s}}^2
  \ge \norm{\by - \yhat}^2,
\]
where the cross term vanishes because $(\by - \yhat)\perp(\yhat-\bm{s})$
(since $\yhat - \bm{s}\in S$), and the Pythagorean theorem applies.
Hence $\yhat$ is the unique minimiser.
\end{proof}


%----------------------------------------------------------------------
\subsection{Characterisation of Orthogonal Projections (Theorem~\ref{thm:proj-char})}
%----------------------------------------------------------------------

\begin{theorem}[Characterisation of Orthogonal Projections]
Let $\bm{P}\in\R^{n\times n}$.  The following are equivalent:
\begin{enumerate}[nosep,label=(\alph*)]
  \item $\bm{P}$ is the matrix of orthogonal projection onto $\Col(\bm{P})$.
  \item $\bm{P}^2 = \bm{P}$ and $\bm{P}^\top = \bm{P}$.
\end{enumerate}
\end{theorem}

\begin{proof}
\textbf{(a) $\Rightarrow$ (b).}
Let $W = \Col(\bm{P})$ and suppose $\bm{P}$ is the orthogonal projection
onto $W$.  For any $\bm{v}\in\R^n$, $\bm{P}\bm{v}\in W$ and
$\bm{v} - \bm{P}\bm{v}\in W^\perp$.

\emph{Idempotency:}
Since $\bm{P}\bm{v}\in W$, projecting it again gives
$\bm{P}(\bm{P}\bm{v}) = \bm{P}\bm{v}$. Hence $\bm{P}^2 = \bm{P}$.

\emph{Symmetry:}
For any $\bm{u}, \bm{v}\in\R^n$,
\begin{align*}
  \ip{\bm{P}\bm{u}}{\bm{v}}
  &= \ip{\bm{P}\bm{u}}{\bm{P}\bm{v} + (\bm{v}-\bm{P}\bm{v})} \\
  &= \ip{\bm{P}\bm{u}}{\bm{P}\bm{v}} + \ip{\bm{P}\bm{u}}{\bm{v}-\bm{P}\bm{v}} \\
  &= \ip{\bm{P}\bm{u}}{\bm{P}\bm{v}},
\end{align*}
because $\bm{P}\bm{u}\in W$ and $\bm{v}-\bm{P}\bm{v}\in W^\perp$.
By the same argument with roles swapped,
$\ip{\bm{u}}{\bm{P}\bm{v}} = \ip{\bm{P}\bm{u}}{\bm{P}\bm{v}}$.
Therefore $\ip{\bm{P}\bm{u}}{\bm{v}} = \ip{\bm{u}}{\bm{P}\bm{v}}$ for all
$\bm{u},\bm{v}$, which means $\bm{P}^\top = \bm{P}$.

\medskip
\noindent\textbf{(b) $\Rightarrow$ (a).}
Assume $\bm{P}^2 = \bm{P}$ and $\bm{P}^\top = \bm{P}$.
Let $W = \Col(\bm{P})$.
We must show: for every $\bm{v}\in\R^n$,
(i) $\bm{P}\bm{v}\in W$, and
(ii) $\bm{v} - \bm{P}\bm{v} \perp W$.

Part (i) is immediate: $\bm{P}\bm{v}$ is a column of $\bm{P}$ (times scalars),
hence in $\Col(\bm{P}) = W$.

For part (ii), let $\bm{w}\in W$ be arbitrary. Then $\bm{w} = \bm{P}\bm{z}$
for some $\bm{z}$.  Compute:
\[
  \ip{\bm{v} - \bm{P}\bm{v}}{\bm{w}}
  = \ip{\bm{v} - \bm{P}\bm{v}}{\bm{P}\bm{z}}
  = (\bm{v} - \bm{P}\bm{v})^\top\bm{P}\bm{z}
  = \bm{v}^\top\bm{P}\bm{z} - \bm{v}^\top\bm{P}^\top\bm{P}\bm{z}.
\]
Using $\bm{P}^\top = \bm{P}$ and then $\bm{P}^2 = \bm{P}$:
\[
  = \bm{v}^\top\bm{P}\bm{z} - \bm{v}^\top\bm{P}^2\bm{z}
  = \bm{v}^\top\bm{P}\bm{z} - \bm{v}^\top\bm{P}\bm{z}
  = 0.
\]
Since $\bm{w}\in W$ was arbitrary, $(\bm{v}-\bm{P}\bm{v})\perp W$.
\end{proof}


%----------------------------------------------------------------------
\subsection{Proof that $h_{ii} \ge 1/n$ when an intercept is present (Remark~\ref{rem:hii-intercept})}
%----------------------------------------------------------------------

When the model includes an intercept, $\bm{1}\in\Col(\bX)$, so
$\bH\bm{1} = \bm{1}$. The $i$-th row of $\bH$ therefore sums to~1:
$\sum_{j=1}^n h_{ij} = 1$.

Since $\bH$ is an orthogonal projection, $h_{ij}^2 \le h_{ii}\cdot h_{jj} \le h_{ii}$
(by Cauchy--Schwarz applied to the rows/columns of $\bH$).  In particular,
\[
  1 = \Bigl(\sum_{j=1}^n h_{ij}\Bigr)^2
  \le n\sum_{j=1}^n h_{ij}^2
  = n\,(\bH^2)_{ii}
  = n\,h_{ii},
\]
where the inequality is Cauchy--Schwarz (applied to the sum $\sum_j h_{ij}\cdot 1$)
and the last step uses $\bH^2 = \bH$.
Therefore $h_{ii}\ge 1/n$.

Without an intercept, it is possible to have $h_{ii} = 0$ (for an observation
$\bm{x}_i$ orthogonal to every column of $\bX$), so the general bound is
$0 \le h_{ii} \le 1$.


\end{document}
